\phantomsection\section*{РЕФЕРАТ}\addcontentsline{toc}{section}{РЕФЕРАТ}

Расчетно-пояснительная записка \pageref{LastPage} с., \total{figure} рис., \total{table} табл., TODO источн., 2 прил. TODO

ТРЕЙСИНГ ПУТИ, ТУМАН, НЕРАВНОМЕРНЫЙ ТУМАН, ОТРАЖЕНИЯ, МОДЕЛЬ ОСВЕЩЕНИЯ, C++, ПУЛ ПОТОКОВ, РЕНДЕР ЧАСТЯМИ

Объект исследования является методы генерации реалистичного изображения сцены и методы визуализации неравномерного тумана в компьютерной графике

Цель работы --- разработать программное обеспечение, использующее алгоритм трассировки пути, для генерации реалистичного изображения сцены,
состоящей из геометрических объектов и неравномерного тумана.

В ходе работы были изучены и реализованы методы объёмного рендеринга, моделирования неоднородных рассеивающих сред,
а также алгоритмы сэмплирования по экспоненциальному распределению и rejection sampling для работы с переменной оптической плотностью.
Реализована программа для фотореалистичной визуализации неоднородного тумана с использованием метода трассировки путей.
Для моделирования пространственной изменчивости плотности среды применён процедурный шум Перлина,
обеспечивающий естественную, облачную структуру тумана.